\snote{synth}

With real data the true parameters are unknown and, in addition, the published estimate $\theta_0$ at which the non-selected parameters are fixed was obtained using the test experiments (\sref{nested}). To measure the prediction error itself, without either limitation, we repeated the external validation with synthetic data generated from a known truth that $\theta_0$ has not seen.

\subsection*{Design}
\begin{enumerate}\itemsep0pt
\item \emph{Truth.} For each seed $s\in\{42,\dots,61\}$, a vector $u\in[-1,1]^p$ with components $u_j\sim\mathcal U(-1,1)$ is drawn (NumPy \texttt{default\_rng}$(s)$) and $\theta^\ast_j=\theta_{0,j}(1+0.2\,u_j)$ for every parameter that can be estimated from the training data (parameters acting only in test conditions stay at $\theta_{0,j}$, since no method could learn them).
\item \emph{Data.} $\tilde y_i=y_i(\theta^\ast)+\epsilon_i$ with $\epsilon_i\sim\mathcal N(0,1)$ in noise-scaled units, i.e.\ Gaussian noise with the standard deviation $\sigma_i$ recorded in the benchmark, drawn from the same generator.
\item \emph{Split and fits.} The same training/test split as for the real data. The reduced model (re-estimating $S\cup\Theta_C$) and the full model (re-estimating all parameters) are fitted to the synthetic training data from $\theta_0$, with the same least-squares method and bounds. The full model was not re-estimated for Chen and Raimundez, which have more than 100 free parameters after removing those that act only in the test conditions (Bachmann has 101 parameters that affect the outputs but 84 free ones, and its full model was re-estimated).
\item \emph{Prediction error.} On the test measurements,
\[ e_{\mathrm{pred}}=\frac{\|y_{\mathrm{test}}(\hat\theta)-y_{\mathrm{test}}(\theta^\ast)\|}{\|y_{\mathrm{test}}(\theta^\ast)-y_{\mathrm{test}}(\theta_0)\|}, \]
the error of the prediction relative to the change of the true response; $e_{\mathrm{pred}}=1$ for $\theta_0$, and $e_{\mathrm{pred}}\le\tau_e=0.436$ is used as the success criterion.
\item \emph{Signal and noise.} When the change from $\theta_0$ to $\theta^\ast$ moves the test outputs less than the noise, the fits are dominated by the noise of the training data and $e_{\mathrm{pred}}$, whose denominator is then small, is unstable. The signal is considered larger than the noise when the test $\chi^2/n$ of $\theta_0$ on the synthetic data is at least 2.
\end{enumerate}
The full-model fits, with many free parameters, can end in different local optima when they are repeated with other library versions, which does not change any qualitative conclusion.

\subsection*{Results}
Over all seeds, 620 problem--seed combinations were planned (31 admissible problems $\times$ 20 seeds). Forty could not be evaluated because Bertozzi and Blasi have no training/test split, eight failed in Isensee because the synthetic truth $\theta^\ast$ could not be simulated, and in two the prediction error is undefined (Chen, seed 51, in which the fitted reduced model could not be simulated on the test conditions, and Perelson, seed 50). The remaining 570 cases, 320 with signal and 250 without, are the denominators below. Among the cases with signal, the reduced model reached $e_{\mathrm{pred}}\le\tau_e$ in 178 of 320 (56\%; between 31\% and 76\% depending on the seed) and predicted at least as well as the full model in 160 of 282 (57\%; 40--73\%). Without signal, the reduced model predicted at the noise level (test $\chi^2/n\le2$) in 197 of 250 cases.

The outcome is stable within each problem (Supplementary Table~\ref{tab:synth20} and main-text Fig.~3b) and divides the benchmark into two groups. The reduced model predicts better than the full model in almost every seed where the full model is over-parameterised for the available data: Weber (20 of 20 seeds), Isensee (12 of 12), Fiedler and Bruno (6 of 6), Brannmark (17 of 18), Zheng (18 of 20) and Elowitz (14 of 17). It predicts worse in almost every seed of the other group, in which the full model, re-estimated on the training data, predicts well (the reason is not established; the full-model profiles of these problems show few identifiable parameters, \sref{profile}): Oliveira (0 of 20), Armistead (1 of 12), Fujita (1 of 18), Okuonghae (2 of 20) and Salazar-Cavazos (3 of 18). Beating the full model does not imply an accurate prediction: in Brannmark the reduced model reaches the threshold in only 1 of 18 seeds, and in Raimundez, Zhao and Schwen it rarely does. Supplementary Table~\ref{tab:synth20seeds} gives the counts for every seed.

{\small
\begin{longtable}{@{}lrrrrrr@{}}
\caption{Synthetic prediction experiment, reduced and full model, 20 seeds (42--61). Seeds: valid seeds; signal: seeds in which the test $\chi^2/n$ of $\theta_0$ on the synthetic data is at least 2. Median $e_{\mathrm{pred}}$ over valid seeds. Last two columns, over the seeds with signal: reduced model with $e_{\mathrm{pred}}\le\tau_e$, and reduced $e_{\mathrm{pred}}$ no larger than that of the full model (--, full model not re-estimated: more than 100 free parameters).}\label{tab:synth20}\\
\toprule
Problem & Seeds & Signal & Reduced $e_{\mathrm{pred}}$ & Full $e_{\mathrm{pred}}$ & Reduced $\le\tau_e$ & Reduced $\le$ full \\
\midrule
\endfirsthead
\multicolumn{7}{@{}l}{\emph{(continued)}}\\
\toprule
Problem & Seeds & Signal & Reduced $e_{\mathrm{pred}}$ & Full $e_{\mathrm{pred}}$ & Reduced $\le\tau_e$ & Reduced $\le$ full \\
\midrule
\endhead
Alkan & 20 & 3 & 0.306 & 0.326 & 3/3 & 1/3 \\
Armistead & 20 & 12 & 0.467 & 0.227 & 7/12 & 1/12 \\
Bachmann & 20 & 2 & 2.85 & 4.65 & 0/2 & 2/2 \\
Bertozzi & \multicolumn{6}{l}{\emph{not evaluable}} \\
Blasi & \multicolumn{6}{l}{\emph{not evaluable}} \\
Boehm & 20 & 4 & 1.51 & 5.65 & 1/4 & 2/4 \\
Borghans & 20 & 11 & 0.648 & 0.494 & 2/11 & 5/11 \\
Brannmark & 20 & 18 & 2.1 & 13.4 & 1/18 & 17/18 \\
Bruno & 20 & 6 & 0.429 & 11 & 4/6 & 6/6 \\
Chen & 19 & 18 & 0.144 & -- & 17/18 & -- \\
Elowitz & 20 & 17 & 0.293 & 0.454 & 11/17 & 14/17 \\
Fiedler & 20 & 6 & 1.15 & 3.51 & 1/6 & 6/6 \\
Fujita & 20 & 18 & 0.345 & 0.183 & 14/18 & 1/18 \\
Giordano & 20 & 0 & 346 & 515 & 0/0 & -- \\
Isensee & 12 & 12 & 0.134 & 0.839 & 9/12 & 12/12 \\
Laske & 20 & 0 & 1.91 & 8.45 & 0/0 & -- \\
Liu & 20 & 1 & 0.863 & 0.913 & 0/1 & 0/1 \\
Lucarelli & 20 & 0 & 1.01 & 3.27 & 0/0 & -- \\
Okuonghae & 20 & 20 & 0.156 & 0.014 & 17/20 & 2/20 \\
Oliveira & 20 & 20 & 0.392 & 0.0111 & 12/20 & 0/20 \\
Perelson & 19 & 1 & 20.9 & 20.9 & 0/1 & 1/1 \\
Rahman & 20 & 19 & 0.118 & 0.372 & 15/19 & 13/19 \\
Raia & 20 & 8 & 0.453 & 0.467 & 7/8 & 5/8 \\
Raimundez & 20 & 20 & 0.571 & -- & 3/20 & -- \\
Salazar-Cavazos & 20 & 18 & 0.334 & 0.213 & 12/18 & 3/18 \\
Schwen & 20 & 9 & 2.63 & 2.53 & 0/9 & 5/9 \\
Smith & 20 & 20 & 0.0954 & 0.179 & 18/20 & 14/20 \\
Sneyd & 20 & 1 & 0.598 & 0.453 & 1/1 & 1/1 \\
Weber & 20 & 20 & 0.359 & 3.77 & 11/20 & 20/20 \\
Zhao & 20 & 16 & 1 & 1 & 0/16 & 11/16 \\
Zheng & 20 & 20 & 0.41 & 0.681 & 12/20 & 18/20 \\
\bottomrule
\end{longtable}

}

{\small
\begin{longtable}{@{}rrrrrr@{}}
\caption{Synthetic prediction experiment, reduced and full model: counts per seed.}\label{tab:synth20seeds}\\
\toprule
Seed & Problems & With signal & Reduced $\le\tau_e$ & Compared with full & Reduced $\le$ full \\
\midrule
\endfirsthead
\multicolumn{6}{@{}l}{\emph{(continued)}}\\
\toprule
Seed & Problems & With signal & Reduced $\le\tau_e$ & Compared with full & Reduced $\le$ full \\
\midrule
\endhead
42 & 28 & 13 & 8 & 11 & 6 \\
43 & 28 & 19 & 12 & 17 & 9 \\
44 & 29 & 16 & 7 & 14 & 7 \\
45 & 28 & 14 & 9 & 12 & 7 \\
46 & 29 & 17 & 9 & 15 & 10 \\
47 & 28 & 14 & 8 & 12 & 6 \\
48 & 29 & 12 & 6 & 11 & 6 \\
49 & 28 & 15 & 10 & 13 & 9 \\
50 & 27 & 17 & 6 & 15 & 11 \\
51 & 28 & 18 & 9 & 17 & 10 \\
52 & 29 & 20 & 11 & 18 & 11 \\
53 & 29 & 17 & 8 & 15 & 6 \\
54 & 29 & 15 & 9 & 13 & 6 \\
55 & 29 & 17 & 10 & 15 & 11 \\
56 & 29 & 16 & 5 & 14 & 6 \\
57 & 29 & 14 & 9 & 12 & 6 \\
58 & 28 & 17 & 13 & 15 & 10 \\
59 & 28 & 17 & 10 & 15 & 8 \\
60 & 29 & 16 & 10 & 14 & 7 \\
61 & 29 & 16 & 9 & 14 & 8 \\
\bottomrule
\end{longtable}

}

\subsection*{Size of the signal and other amplitudes}
Because the outputs are scaled by the recorded noise, the change of the true response, $y(\theta^\ast)-y(\theta_0)$, can be measured directly in noise units, without the noisy data: its root mean square per test measurement (\texttt{sintetico\_semillas.py} with \texttt{SINT\_SOLO\_SENAL=1}). With the amplitude of the experiment ($\theta^\ast=\theta_0(1+0.2u)$) its median over all problems and seeds is $1.25\sigma$ (interquartile range 0.50--2.84). The success of the reduced model grows with this signal: $e_{\mathrm{pred}}\le\tau_e$ in 1\% of the cases with a signal below $0.3\sigma$, 24\% between $0.3\sigma$ and $\sigma$, 44\% between $\sigma$ and $3\sigma$ and 75\% above $3\sigma$. Below about one noise unit $e_{\mathrm{pred}}$ therefore reflects mainly the noise of the training data. The criterion used above to define the cases with signal (test $\chi^2/n$ of $\theta_0$ at least 2), which uses the noisy data, agrees closely with this direct measure: it classifies as cases with signal 93--100\% of those whose true signal exceeds $\sigma$ and only 1--8\% of the others.

To evaluate the dependence on the distance between the truth and $\theta_0$, the experiment was repeated with amplitudes $\pm5\%$ and $\pm50\%$ (seeds 42--51; \texttt{SINT\_AMPLITUD}) and compared with $\pm20\%$ on the same seeds. The median signal is $0.31\sigma$, $1.25\sigma$ and $2.63\sigma$ at $\pm5\%$, $\pm20\%$ and $\pm50\%$, so the number of cases with signal grows from 64 to 155 and 219. Among them, the reduced model reached $e_{\mathrm{pred}}\le\tau_e$ in 64\%, 54\% and 51\% of the cases, and predicted at least as well as the full model in 67\% (42 of 63), 59\% (81 of 137) and 47\% (95 of 201). The reduction is local: it predicts best when the truth is close to the reference $\theta_0$ around which the subset was selected, and its advantage over the full model disappears for changes of $\pm50\%$, although it still reaches the threshold in about half of the cases.
